% ============================================================
% Chapter 13: Cyber Resilience, Governance and Future Trends
% Language: Greek — edit freely; regenerate from src/content if preferred.
% ============================================================
\chapter{Κυβερνοανθεκτικότητα, Διακυβέρνηση και Μελλοντικές Τάσεις}\label{ch:13}
\chaptermeta{Αντίγραφα ασφαλείας και ανάκαμψη · BCP, RPO και RTO · Διαχείριση κινδύνου · NIST CSF 2.0, ISO/IEC 27001, CIS · ΓΚΠΔ και NIS2 · Δείκτες · Ασφάλεια OT · Μηδενική εμπιστοσύνη και νέφος}{Επίπεδο: Προχωρημένο επίπεδο · Στρατηγική και διακυβέρνηση}{Χρόνος μελέτης: 10–12 ώρες}

Το τελευταίο κεφάλαιο επιστρέφει στο ερώτημα που τέθηκε στο Κεφάλαιο 1: \emph{μπορεί ο οργανισμός να συνεχίσει να λειτουργεί και να ανακάμψει όταν τα πράγματα πάνε στραβά;} Τα τεχνικά μέτρα είναι αναγκαία αλλά όχι επαρκή. Η ανθεκτικότητα εξαρτάται επίσης από αντίγραφα ασφαλείας που λειτουργούν όταν χρειάζονται, από σχέδια που έχουν δοκιμαστεί, από αποφάσεις για τον κίνδυνο που λαμβάνονται συνειδητά από υπόλογα πρόσωπα και από τη συμμόρφωση με ένα διαρκώς διευρυνόμενο νομικό πλαίσιο.

Το κεφάλαιο καλύπτει πρώτα τη μηχανική ανθεκτικότητας, τις αρχιτεκτονικές αντιγράφων ασφαλείας και την επιχειρησιακή συνέχεια. Στη συνέχεια στρέφεται στη διακυβέρνηση: τη διαχείριση κινδύνου στην πράξη, τα κυριότερα πλαίσια και κανονιστικά κείμενα και τους δείκτες που δείχνουν στη διοίκηση αν η ασφάλεια βελτιώνεται. Ολοκληρώνεται εξετάζοντας περιβάλλοντα και τάσεις που θα διαμορφώσουν τα επόμενα χρόνια —την επιχειρησιακή τεχνολογία και τις κρίσιμες υποδομές, την αρχιτεκτονική μηδενικής εμπιστοσύνης και την ασφάλεια του νέφους και της εφοδιαστικής αλυσίδας.

\begin{outcomesbox}{Μαθησιακά αποτελέσματα}
Μετά τη μελέτη του κεφαλαίου θα πρέπει να μπορείτε:
\begin{itemize}
  \item Να σχεδιάζετε στρατηγική αντιγράφων ασφαλείας που αντέχει σε λυτρισμικό και να ορίζετε RPO και RTO για κρίσιμες υπηρεσίες.
  \item Να καταρτίζετε μητρώο κινδύνων και να επιλέγετε κατάλληλες επιλογές αντιμετώπισης.
  \item Να συγκρίνετε τα NIST CSF 2.0, ISO/IEC 27001 και CIS Controls και να συνοψίζετε τις υποχρεώσεις του ΓΚΠΔ και της NIS2.
  \item Να διακρίνετε τους ουσιαστικούς δείκτες ασφάλειας από τους δείκτες «βιτρίνας».
  \item Να εξηγείτε τους ιδιαίτερους περιορισμούς της ασφάλειας OT, καθώς και τις αρχές της μηδενικής εμπιστοσύνης και της κοινής ευθύνης στο νέφος.
\end{itemize}
\end{outcomesbox}

\section{Μηχανική ανθεκτικότητας και αρχιτεκτονικές αντιγράφων ασφαλείας}\label{sec:13.1}
Η \textbf{μηχανική ανθεκτικότητας} αποδέχεται ότι ορισμένες επιθέσεις και αστοχίες θα επιτύχουν και σχεδιάζει τα συστήματα ώστε να υποβαθμίζονται ομαλά αντί να καταρρέουν. Τα εργαλεία της περιλαμβάνουν πλεονασμό χωρίς κοινά μεμονωμένα σημεία αστοχίας, δυνατότητα απομόνωσης των στοιχείων που έχουν υποστεί ζημιά και προσχεδιασμένους τρόπους περιορισμένης λειτουργίας —για παράδειγμα, ένα νοσοκομείο που μπορεί να συνεχίσει να δέχεται ασθενείς σε χαρτί όταν το ηλεκτρονικό του σύστημα δεν είναι διαθέσιμο.

Τα αντίγραφα ασφαλείας αποτελούν το θεμέλιο της ανάκαμψης, και το σύγχρονο λυτρισμικό τα στοχεύει σκόπιμα. Ο κλασικός \textbf{κανόνας 3-2-1} —τρία αντίγραφα των δεδομένων, σε δύο διαφορετικούς τύπους μέσων, ένα από τα οποία εκτός εγκατάστασης— έχει γι' αυτό επεκταθεί σε \textbf{3-2-1-1-0}: τουλάχιστον ένα αντίγραφο πρέπει να είναι \textbf{εκτός σύνδεσης ή αμετάβλητο} (αποθήκευση εφάπαξ εγγραφής ή κλείδωμα αντικειμένων που ούτε ένας διαχειριστής δεν μπορεί να διαγράψει κατά την περίοδο διατήρησης), και οι επαναφορές πρέπει να δοκιμάζονται, ώστε να υπάρχουν \textbf{μηδέν} ανεπιβεβαίωτα σφάλματα. Τα συστήματα αντιγράφων ασφαλείας χρειάζονται δικά τους, ξεχωριστά διαπιστευτήρια και MFA, διότι διαφορετικά ένας επιτιθέμενος με δικαιώματα διαχειριστή τομέα απλώς θα τα διαγράψει πριν ξεκινήσει την κρυπτογράφηση.

\section{Επιχειρησιακή συνέχεια και ανάκαμψη από καταστροφές: RPO και RTO}\label{sec:13.2}
Ο \textbf{σχεδιασμός επιχειρησιακής συνέχειας} (Business Continuity Planning, BCP) διασφαλίζει ότι οι κρίσιμες επιχειρησιακές λειτουργίες συνεχίζονται κατά τη διάρκεια μιας διαταραχής, ενώ η \textbf{ανάκαμψη από καταστροφές} (Disaster Recovery, DR) εστιάζει στην αποκατάσταση των συστημάτων πληροφορικής που τις υποστηρίζουν. Και οι δύο ξεκινούν με μια \textbf{ανάλυση επιχειρησιακών επιπτώσεων}, η οποία εντοπίζει τις κρίσιμες διεργασίες και τις συνέπειες της διακοπής τους σε βάθος χρόνου. Δύο παράμετροι καθοδηγούν στη συνέχεια τον τεχνικό σχεδιασμό. Ο \textbf{στόχος σημείου ανάκαμψης} (Recovery Point Objective, RPO) είναι η μέγιστη αποδεκτή απώλεια δεδομένων, μετρούμενη σε χρόνο: RPO μίας ώρας απαιτεί αντίγραφα ή αναπαραγωγή τουλάχιστον ανά ώρα. Ο \textbf{στόχος χρόνου ανάκαμψης} (Recovery Time Objective, RTO) είναι ο μέγιστος αποδεκτός χρόνος αποκατάστασης της υπηρεσίας. Οι αυστηρότεροι στόχοι κοστίζουν περισσότερο και πρέπει, επομένως, να δικαιολογούνται από τις επιχειρησιακές επιπτώσεις.

Τα σχέδια είναι αξιόπιστα μόνο εφόσον δοκιμάζονται. Τα \textbf{εγχειρίδια ανάκαμψης} (DR playbooks) περιγράφουν βήμα προς βήμα πώς αποκαθίσταται κάθε κρίσιμο σύστημα, με ποια σειρά (οι υπηρεσίες ταυτότητας και η δικτύωση προηγούνται συνήθως) και ποιος είναι υπεύθυνος. Επικυρώνονται μέσω ασκήσεων επί χάρτου (tabletop), μερικών δοκιμών επαναφοράς και πλήρων ασκήσεων μετάπτωσης, και οι μετρούμενοι χρόνοι ανάκαμψης συγκρίνονται με τον RTO. Ένα σχέδιο που δεν έχει δοκιμαστεί ποτέ πρέπει να θεωρείται ότι δεν λειτουργεί.

\section{Διαχείριση επιχειρησιακού κινδύνου στην πράξη}\label{sec:13.3}
Το Κεφάλαιο 1 όρισε τον κίνδυνο· η διακυβέρνηση μετατρέπει αυτόν τον ορισμό σε διαδικασία διαχείρισης. Ένα \textbf{μητρώο κινδύνων} καταγράφει κάθε κίνδυνο με σαφή διατύπωση (\emph{απειλή} που εκμεταλλεύεται \emph{ευπάθεια} προκαλώντας \emph{επίπτωση} σε \emph{περιουσιακό στοιχείο}), υπεύθυνο, υφιστάμενα μέτρα, εκτιμήσεις πιθανότητας και επίπτωσης, την επιλεγμένη αντιμετώπιση και ημερομηνία επανεξέτασης. Ένας \textbf{πίνακας κινδύνων} —συνήθως πέντε επί πέντε— βοηθά στην οπτικοποίηση και σύγκριση των κινδύνων, αν και οι ποιοτικές κλίμακες πρέπει να ορίζονται προσεκτικά, ώστε το «πιθανό» να σημαίνει το ίδιο για όλους. Ποσοτικές μέθοδοι όπως η FAIR εκφράζουν τον κίνδυνο σε χρηματικούς όρους και μπορούν να υποστηρίξουν αποφάσεις επενδύσεων.

Κάθε κίνδυνος πρέπει να λαμβάνει ρητή \textbf{απόφαση αντιμετώπισης}: \emph{μείωση} με πρόσθετα μέτρα, \emph{μεταφορά} μέσω ασφάλισης ή συμβάσεων, \emph{αποφυγή} με διακοπή της δραστηριότητας ή επίσημη \emph{αποδοχή}. Η αποδοχή είναι θεμιτή μόνο όταν αποφασίζεται από πρόσωπο με την αρμοδιότητα να φέρει τις συνέπειες, τεκμηριώνεται και επανεξετάζεται περιοδικά. Οι τεχνικές βαθμολογίες, όπως το CVSS (Κεφάλαιο 10), αποτελούν δεδομένα εισόδου αυτής της διαδικασίας και όχι υποκατάστατό της.

\section{Πλαίσια και ρυθμίσεις: NIST CSF 2.0, ISO/IEC 27001, CIS, ΓΚΠΔ και NIS2}\label{sec:13.4}
Τα πλαίσια δίνουν δομή σε ένα πρόγραμμα ασφάλειας. Το \textbf{NIST Cybersecurity Framework 2.0} (2024) οργανώνει τα επιδιωκόμενα αποτελέσματα σε έξι λειτουργίες: \emph{Διακυβέρνηση}, \emph{Αναγνώριση}, \emph{Προστασία}, \emph{Ανίχνευση}, \emph{Απόκριση} και \emph{Ανάκαμψη}· η νέα λειτουργία της Διακυβέρνησης υπογραμμίζει ότι η κυβερνοασφάλεια είναι ευθύνη της ηγεσίας. Το \textbf{ISO/IEC 27001} ορίζει τις απαιτήσεις ενός \emph{συστήματος διαχείρισης ασφάλειας πληροφοριών} (ΣΔΑΠ/ISMS) —ενός κύκλου αξιολόγησης κινδύνου, επιλογής μέτρων, παρακολούθησης και βελτίωσης— και αποτελεί τη βάση για ανεξάρτητη πιστοποίηση· το Παράρτημα Α περιλαμβάνει 93 μέτρα αναφοράς. Τα \textbf{CIS Critical Security Controls} είναι ένα ιεραρχημένο σύνολο συγκεκριμένων τεχνικών μέτρων, ομαδοποιημένων σε ομάδες υλοποίησης για οργανισμούς διαφορετικής ωριμότητας.

Το κανονιστικό πλαίσιο καθιστά ορισμένες πρακτικές υποχρεωτικές. Ο \textbf{Γενικός Κανονισμός για την Προστασία Δεδομένων} (ΓΚΠΔ/GDPR) της ΕΕ απαιτεί κατάλληλη ασφάλεια για τα προσωπικά δεδομένα, προστασία δεδομένων ήδη από τον σχεδιασμό και εξ ορισμού, καθώς και γνωστοποίηση των παραβιάσεων προσωπικών δεδομένων στην εποπτική αρχή εντός 72 ωρών. Η \textbf{Οδηγία NIS2} επεκτείνει τις υποχρεώσεις κυβερνοασφάλειας σε ευρύ φάσμα βασικών και σημαντικών οντοτήτων —ενέργεια, υγεία, μεταφορές, ψηφιακές υποδομές κ.ά.— απαιτώντας μέτρα διαχείρισης κινδύνου, ασφάλεια της εφοδιαστικής αλυσίδας, αναφορά περιστατικών (έγκαιρη προειδοποίηση εντός 24 ωρών) και προσωπική λογοδοσία των διοικητικών οργάνων. Στις Ηνωμένες Πολιτείες, ο νόμος \textbf{HIPAA} επιβάλλει ανάλογες διασφαλίσεις για τις πληροφορίες υγείας.

\begin{table}[htbp]
  \centering\small
  \caption{Πλαίσια και κανονιστικά κείμενα με μια ματιά}
  \begin{tabularx}{\linewidth}{>{\raggedright\arraybackslash}X >{\raggedright\arraybackslash}X >{\raggedright\arraybackslash}X >{\raggedright\arraybackslash}X}
    \toprule
    \textbf{Κείμενο} & \textbf{Φύση} & \textbf{Κύρια εστίαση} & \textbf{Τυπική χρήση} \\
    \midrule
    NIST CSF 2.0 & Προαιρετικό πλαίσιο & Αποτελέσματα σε έξι λειτουργίες & Δομή προγράμματος, αποτύπωση ωριμότητας \\
    ISO/IEC 27001 & Πιστοποιήσιμο πρότυπο & Σύστημα διαχείρισης (ΣΔΑΠ) & Πιστοποίηση, διαβεβαίωση πελατών \\
    CIS Controls v8 & Ιεραρχημένο σύνολο μέτρων & Συγκεκριμένα τεχνικά μέτρα & Οδικός χάρτης υλοποίησης \\
    ΓΚΠΔ & Κανονισμός ΕΕ (δεσμευτικός) & Προστασία προσωπικών δεδομένων & Νομική συμμόρφωση, γνωστοποίηση παραβιάσεων \\
    NIS2 & Οδηγία ΕΕ (δεσμευτική) & Ασφάλεια βασικών/σημαντικών οντοτήτων & Διαχείριση κινδύνου, αναφορά περιστατικών \\
    \bottomrule
  \end{tabularx}
\end{table}
\section{Δείκτες ασφάλειας: ουσιαστική πληροφορία έναντι βιτρίνας}\label{sec:13.5}
Η διοίκηση χρειάζεται τεκμήρια ότι η επένδυση στην ασφάλεια μειώνει τον κίνδυνο. Πολλοί συνήθως αναφερόμενοι αριθμοί είναι \textbf{δείκτες βιτρίνας}: τα «εκατομμύρια επιθέσεων που αποκλείστηκαν» μετρούν κυρίως τον αυτοματοποιημένο «θόρυβο» του διαδικτύου, ενώ ο «αριθμός των εργαλείων που έχουν εγκατασταθεί» δεν λέει τίποτα για την αποτελεσματικότητά τους. Οι \textbf{ουσιαστικοί δείκτες} συνδέονται με αποφάσεις και αποτελέσματα: το ποσοστό των κρίσιμων περιουσιακών στοιχείων που καλύπτονται από EDR και καταγραφή· ο χρόνος εφαρμογής ενημερώσεων για ευπάθειες σε συστήματα εκτεθειμένα στο διαδίκτυο που είναι γνωστό ότι αποτελούν αντικείμενο εκμετάλλευσης· οι MTTD και MTTR (Κεφάλαιο 11)· το ποσοστό αναφοράς μηνυμάτων ψαρέματος (Κεφάλαιο 5)· το ποσοστό των προνομιούχων λογαριασμών που προστατεύονται με MFA ανθεκτικό στο ψάρεμα· και το ποσοστό επιτυχίας και η διάρκεια των δοκιμών επαναφοράς αντιγράφων ασφαλείας σε σύγκριση με τον RTO.

Οι καλοί δείκτες έχουν καθορισμένη πηγή δεδομένων, στόχο και τάση στον χρόνο, και ο καθένας πρέπει να υποδεικνύει μια ενέργεια αν κινηθεί προς τη λάθος κατεύθυνση. Η συνεπής αναφορά λίγων τέτοιων δεικτών είναι πολύ πιο χρήσιμη για ένα διοικητικό συμβούλιο από έναν πίνακα με εκατοντάδες ανεξήγητα νούμερα.

\section{Κρίσιμες υποδομές και επιχειρησιακή τεχνολογία}\label{sec:13.6}
Η \textbf{επιχειρησιακή τεχνολογία} (Operational Technology, OT) —βιομηχανικά συστήματα ελέγχου (ICS), συστήματα SCADA και προγραμματιζόμενοι λογικοί ελεγκτές που λειτουργούν δίκτυα ηλεκτρικής ενέργειας, εγκαταστάσεις επεξεργασίας νερού, παραγωγικές μονάδες και ιατροτεχνολογικά προϊόντα— έχει προτεραιότητες διαφορετικές από την πληροφορική γραφείου. Η ασφάλεια των ανθρώπων και η διαθεσιμότητα προηγούνται, τα συστήματα μπορεί να λειτουργούν για είκοσι χρόνια ή και περισσότερο, και πολλά βιομηχανικά πρωτόκολλα, όπως το Modbus, σχεδιάστηκαν χωρίς έλεγχο ταυτότητας. Η εφαρμογή ενημερώσεων μπορεί να απαιτεί διακοπή μιας γραμμής παραγωγής, ενώ μια επιθετική δικτυακή σάρωση μπορεί να θέσει εκτός λειτουργίας έναν ευαίσθητο ελεγκτή. Όπως έδειξαν το Stuxnet και μεταγενέστερες επιθέσεις σε ενεργειακά δίκτυα, οι συνέπειες μπορεί να είναι φυσικές.

Η ασφάλεια OT βασίζεται επομένως σε μεγάλο βαθμό στην αρχιτεκτονική. Το \textbf{μοντέλο Purdue} διαιρεί το περιβάλλον σε επίπεδα, από τις φυσικές διεργασίες (επίπεδο 0) και τους ελεγκτές (επίπεδο 1) έως την επιχειρησιακή πληροφορική (επίπεδα 4–5). Το πρότυπο \textbf{ISA/IEC 62443} ομαδοποιεί τα στοιχεία σε \emph{ζώνες} και επιτρέπει την επικοινωνία μόνο μέσω ελεγχόμενων \emph{αγωγών} (conduits), με μια βιομηχανική αποστρατιωτικοποιημένη ζώνη μεταξύ IT και OT. Η παθητική παρακολούθηση του δικτύου, σχεδιασμένη για βιομηχανικά πρωτόκολλα, παρέχει ορατότητα χωρίς να διαταράσσει τις ευαίσθητες συσκευές, ενώ ο αυστηρός έλεγχος της απομακρυσμένης πρόσβασης των προμηθευτών κλείνει ένα από τα σημεία εισόδου που γίνονται συχνότερα αντικείμενο κατάχρησης.

\section{Μηδενική εμπιστοσύνη, ασφάλεια νέφους και αναδυόμενες απειλές}\label{sec:13.7}
Η \textbf{αρχιτεκτονική μηδενικής εμπιστοσύνης} (Zero Trust Architecture, NIST SP 800-207) εγκαταλείπει την ιδέα ότι οτιδήποτε βρίσκεται μέσα στην περίμετρο του δικτύου είναι αξιόπιστο. Κάθε αίτημα πρόσβασης αξιολογείται ρητά, με βάση την ταυτότητα του χρήστη, την κατάσταση της συσκευής και την ευαισθησία του πόρου· η πρόσβαση παραχωρείται με ελάχιστο προνόμιο, για περιορισμένη συνεδρία· και ο σχεδιασμός προϋποθέτει ότι μια παραβίαση έχει ήδη συμβεί. Ένα \emph{σημείο λήψης απόφασης πολιτικής} αξιολογεί κάθε αίτημα και ένα \emph{σημείο επιβολής πολιτικής} μπροστά από κάθε πόρο εφαρμόζει την απόφαση. Η μηδενική εμπιστοσύνη είναι πορεία και όχι προϊόν: οι οργανισμοί συνήθως προχωρούν από την ισχυρή ταυτότητα και το MFA, στη συμμόρφωση των συσκευών και στη μικροκατάτμηση και, τέλος, σε συνεχείς αποφάσεις πρόσβασης προσαρμοσμένες στον κίνδυνο.

Στο \textbf{νέφος}, η ασφάλεια ακολουθεί ένα \textbf{μοντέλο κοινής ευθύνης}: ο πάροχος προστατεύει την υποκείμενη υποδομή, ενώ ο πελάτης παραμένει υπεύθυνος για τις ταυτότητες, τη διαμόρφωση και τα δεδομένα. Οι περισσότερες παραβιάσεις στο νέφος οφείλονται σε λανθασμένες ρυθμίσεις από την πλευρά του πελάτη —δημόσια αναγνώσιμους χώρους αποθήκευσης, υπερβολικά δικαιώματα, εκτεθειμένα κλειδιά— τις οποίες έχουν σχεδιαστεί να εντοπίζουν τα εργαλεία \emph{διαχείρισης της στάσης ασφάλειας στο νέφος} (CSPM). Τα κοντέινερ προσθέτουν τις δικές τους απαιτήσεις: ελάχιστα και ελεγμένα είδωλα, εκτέλεση χωρίς δικαιώματα root και πολιτικές δικτύου Kubernetes με άρνηση εξ ορισμού. Κοιτάζοντας μπροστά, οι αμυνόμενοι πρέπει να προετοιμαστούν για τη μετάβαση στη \textbf{μετακβαντική κρυπτογραφία}, για το ηλεκτρονικό ψάρεμα με υποβοήθηση τεχνητής νοημοσύνης και τις απάτες με deepfakes, καθώς και για τη συνέχιση των επιθέσεων στην εφοδιαστική αλυσίδα λογισμικού. Οι αρχές αυτού του βιβλίου —ελάχιστο προνόμιο, άμυνα σε βάθος, επαλήθευση και ανθεκτικότητα— παραμένουν ο πιο αξιόπιστος οδηγός μέσα σε αυτές τις αλλαγές.

\section*{Βασικοί όροι}
\begin{description}[style=nextline,leftmargin=1.2em]
  \item[3-2-1-1-0] Κανόνας αντιγράφων: 3 αντίγραφα, 2 μέσα, 1 εκτός εγκατάστασης, 1 εκτός σύνδεσης/αμετάβλητο, 0 ανεπιβεβαίωτα σφάλματα επαναφοράς.
  \item[RPO / RTO] Μέγιστη αποδεκτή απώλεια δεδομένων (RPO) και μέγιστος αποδεκτός χρόνος διακοπής (RTO).
  \item[Μητρώο κινδύνων] Δομημένη καταγραφή κινδύνων, υπευθύνων, μέτρων, εκτιμήσεων και αποφάσεων αντιμετώπισης.
  \item[ΣΔΑΠ (ISMS)] Σύστημα διαχείρισης ασφάλειας πληροφοριών, όπως ορίζεται στο ISO/IEC 27001.
  \item[Μοντέλο Purdue] Πολυεπίπεδη αρχιτεκτονική αναφοράς που διαχωρίζει τα επίπεδα βιομηχανικού ελέγχου από την επιχειρησιακή πληροφορική.
  \item[Κοινή ευθύνη] Κατανομή των καθηκόντων ασφάλειας μεταξύ ενός παρόχου νέφους και του πελάτη του.
\end{description}

\section*{Σύνοψη κεφαλαίου}
\begin{itemize}
  \item Η ανθεκτικότητα προϋποθέτει την αστοχία: τα αμετάβλητα, δοκιμασμένα αντίγραφα ασφαλείας με ξεχωριστά διαπιστευτήρια αποτελούν τον πυρήνα της ανάκαμψης από λυτρισμικό.
  \item Η επιχειρησιακή συνέχεια και η ανάκαμψη από καταστροφές καθοδηγούνται από τις επιχειρησιακές επιπτώσεις, εκφράζονται μέσω RPO και RTO και είναι αξιόπιστες μόνο όταν δοκιμάζονται.
  \item Η διαχείριση κινδύνου καταγράφει, αξιολογεί και αντιμετωπίζει ρητά κάθε κίνδυνο, με την αποδοχή να αποφασίζεται από υπόλογους υπευθύνους.
  \item Τα NIST CSF 2.0, ISO/IEC 27001 και CIS Controls δομούν τα προγράμματα ασφάλειας· ο ΓΚΠΔ και η NIS2 καθιστούν υποχρεωτικές βασικές πρακτικές.
  \item Η ασφάλεια OT, η μηδενική εμπιστοσύνη και η κοινή ευθύνη στο νέφος επεκτείνουν τις βασικές αρχές του βιβλίου σε νέα περιβάλλοντα.
\end{itemize}

\section*{Ερωτήσεις ανασκόπησης}
\begin{enumerate}
  \item Γιατί ο κανόνας 3-2-1 από μόνος του δεν εγγυάται την ανάκαμψη από λυτρισμικό; Τι προσθέτουν το επιπλέον «1» και το «0»;
  \item Ένα ηλεκτρονικό κατάστημα μπορεί να ανεχθεί απώλεια παραγγελιών 15 λεπτών και διακοπή λειτουργίας 2 ωρών. Μετατρέψτε αυτά σε RPO/RTO και προτείνετε τεχνικό σχεδιασμό.
  \item Ποιος πρέπει να επιτρέπεται να αποδεχθεί έναν υψηλό υπολειπόμενο κίνδυνο και γιατί η αποδοχή πρέπει να τεκμηριώνεται;
  \item Εξηγήστε γιατί η ενεργητική σάρωση ευπαθειών, συνήθης στην πληροφορική, μπορεί να είναι ακατάλληλη σε περιβάλλον OT.
\end{enumerate}
